\section{Introduction}\label{sec:introduction}

The Euclidean isoperimetric inequality bounds the volume enclosed by a
boundary of prescribed area. Its current-theoretic form asks for a filling
of a cycle and permits singular supports, integer multiplicities, and
arbitrary codimension. We prove that the sharp Euclidean filling bound
holds in every proper CAT(0) space, in all cycle dimensions at least two.

Write $\I_j(X)$ for the integral $j$-currents of Ambrosio and
Kirchheim~\cite{AK2000}, $\partial$ for current boundary, and $\M$
for mass. An integral $n$-cycle is a current $T\in\I_n(X)$ with
$\partial T=0$; an integral filling is a current $S\in\I_{n+1}(X)$
with $\partial S=T$. A metric space is \emph{proper} if its closed
bounded balls are compact. A CAT(0) space is a geodesic metric space in which the distance
between two points on a geodesic triangle is at most the distance between
the corresponding points of its Euclidean comparison triangle.

Let $\omega_j$ be the volume of the unit ball in $\R^j$ and put
\begin{equation}\label{eq:constants}
 s_n=(n+1)\omega_{n+1}=\operatorname{area}(\mathbb S^n),
 \qquad c_n=\frac{1}{(n+1)s_n^{1/n}}.
\end{equation}
Here and throughout the filling argument, $n$ denotes the cycle
dimension.

\begin{theorem}[Sharp integral filling]\label{thm:main}
Let $X$ be a proper CAT(0) space and $n\ge2$ an integer.
Every compactly supported $T\in\I_n(X)$ with $\partial T=0$ admits a
compactly supported $S\in\I_{n+1}(X)$ such that
\begin{equation}\label{eq:main}
 \partial S=T,\qquad
 \M(S)\le c_n\M(T)^{(n+1)/n}
 =\frac{\M(T)^{(n+1)/n}}
 {(n+1)^{(n+1)/n}\omega_{n+1}^{1/n}}.
\end{equation}
There is no restriction on the dimension of $X$ or on the integer
multiplicities of the currents. The coefficient is optimal.
\end{theorem}

For the first two dimensions, $s_2=4\pi$ and $s_3=2\pi^2$, so
\[
 c_2=\frac1{6\sqrt\pi},\qquad
 c_3=(128\pi^2)^{-1/3}.
\]
Euclidean spheres attain the coefficient, as verified in
Section~\ref{sec:classical}.

\subsection{The Cartan--Hadamard consequence}

A \emph{Cartan--Hadamard manifold} is a complete, simply connected
Riemannian manifold of nonpositive sectional curvature. The Euclidean
Cartan--Hadamard conjecture asks whether domains in such a manifold
satisfy the same volume--perimeter bound as Euclidean domains.
Theorem~\ref{thm:main}, applied to boundary currents, gives the
following positive answer in dimensions at least three.

\begin{corollary}\label{cor:cartan-hadamard}
Let $M$ be a Cartan--Hadamard manifold of dimension $d\ge3$.
For every bounded domain $\Omega\subset M$ with smooth boundary,
\begin{equation}\label{eq:classical}
 \operatorname{area}(\partial\Omega)
 \ge d\omega_d^{1/d}\vol(\Omega)^{(d-1)/d}.
\end{equation}
The same inequality holds for relatively compact sets of finite
perimeter, with ambient perimeter in place of boundary area.
\end{corollary}

The reduction uses the uniqueness of a compactly supported
top-dimensional filling of a prescribed boundary. Its proof, and the
classical two-dimensional case, appear in Section~\ref{sec:classical}.
Together they give the Euclidean inequality for bounded smooth domains
in every ambient dimension $d\ge2$.

The smooth domain problem has a long independent history. The surface
inequality goes back to Weil and Beckenbach--Rad\'o
\cite{Weil1926,BeckenbachRado1933}; Kleiner proved the
three-dimensional comparison~\cite{Kleiner1992}, and Croke proved
the four-dimensional Euclidean inequality~\cite{Croke1984}.
Ghomi and Spruck showed that a corresponding total-curvature bound
implies the domain inequality in every dimension
\cite[Theorem~7.1]{GhomiSpruck2022}.

Recent results extend these comparisons in several directions. Chen,
Ghomi and Wang prove the Euclidean comparison in dimensions three
through nine~\cite[Theorem~1.1]{ChenGhomiWang2026HigherDimensions}.
Wheeler proves it in all dimensions under a controlled variation
condition on the negative-curvature scale
\cite[Theorem~1.1 and Definition~1.2]{Wheeler2026}.
Agnoletto, Da Silva, Nardulli and Resende prove small-volume comparison
with the constant-curvature model under a nonpositive sectional-curvature
upper bound and a Ricci lower bound,
and reduce unrestricted-volume comparison to equality rigidity in a
class of Alexandrov spaces
\cite[Theorems~1.1 and~1.3; Assumption~1]{AgnolettoDaSilvaNardulliResende2026}.
Our domain corollary follows from the integral filling theorem and
top-dimensional constancy; it requires no separate smooth comparison
argument or additional curvature hypothesis.

\subsection{The filling problem and its predecessors}

Almgren established the sharp isoperimetric inequality for integral
currents in Euclidean space in arbitrary dimension and
codimension~\cite{Almgren1986}. In a general metric space, a comparable
statement first requires notions of orientation, integer multiplicity,
mass, and boundary that do not depend on an ambient smooth structure.
Ambrosio and Kirchheim supplied this theory, including rectifiable
representation, slicing, and compactness~\cite{AK2000}.
Kirchheim's metric differentiation theorem~\cite{Kirchheim1994}
provides the local normed geometry of its rectifiable charts. In CAT(0)
spaces these chart norms are Euclidean; we give the short comparison
argument and the exact mass normalization in Section~\ref{sec:currents}.

For Lipschitz loops in complete CAT(0) spaces, Reshetnyak's majorization
theorem already gives a Lipschitz disc whose parametrized Hausdorff area
is at most the squared loop length divided by $4\pi$; see
\cite[Section~3.2, Lemma~3.2]{LytchakWenger2018}.
Wenger proved isoperimetric inequalities with the Euclidean exponent
in complete spaces with a convex geodesic bicombing, including CAT(0)
spaces~\cite{Wenger2005Iso}. His constant depends on the cycle
dimension. He also established the relation between weak convergence
and filling convergence used to obtain variational extremizers
here~\cite{Wenger2007Flat}. Theorem~\ref{thm:main} identifies the optimal
coefficient. For two-cycles in smooth Cartan--Hadamard manifolds,
Schulze proved the sharp constant-curvature model comparison in every
codimension, together with equality rigidity
\cite[Theorems~1.2--1.3]{Schulze2020}. The present theorem extends the
Euclidean filling bound to all cycle dimensions $n\ge2$ and singular CAT(0)
ambient spaces.

Recent work identifies geometric hypotheses under which fillings
have smaller growth exponents at large mass. In complete CAT(0)
spaces of finite asymptotic Nagata dimension, Lang, Stadler and Urech
obtain almost-linear bounds for compact integral $k$-cycles, $k\ge2$,
when the asymptotic rank is at most two
\cite[Theorem~1.1]{LangStadlerUrech2026}. With the same Nagata dimension
hypothesis and finite asymptotic rank $k_0$, Peteranderl obtains linear
bounds for integral $k$-cycles with $k\ge\max\{k_0,1\}$, and compact
fillings for compact inputs~\cite[Theorem~1.2]{Peteranderl2026}.
The constants depend on the space. The estimate here fixes the universal
Euclidean coefficient at every mass in an arbitrary proper CAT(0) space.

The problem also belongs to the broader study of filling geometry
developed by Gromov~\cite{Gromov1983}. His sharp shrinking and filling
formulation for CAT(0) targets asks for volume-controlled extensions of
maps on specified domains~\cite[Section~2.3]{Gromov2014}.
Theorem~\ref{thm:main} concerns integral-current fillings: it does not
prescribe the topology or parametrization of the filling. Those
additional extension requirements are not supplied by a mass and
boundary estimate alone.

\subsection{Main ideas and proof route}

The extremizer/curvature-estimate strategy is motivated by Almgren
\cite[Section~0, p.~454]{Almgren1986}, and the two-dimensional
curvature estimate also by Schulze \cite[Theorem~1.4]{Schulze2020}.
The singular CAT(0) radial implementation is established here; these
references are methodological motivation, not imports of a smooth
curvature theorem into the singular setting.

The proof proceeds by induction on $n$, with its two-dimensional case
proved directly. A closed convex ball containing the given compact
support reduces the problem to a compact CAT(0) space $Y$. For an
integral cycle $B$ in $Y$, let $V(B)$ be its least integral filling mass.
If the sharp bound fails, compactness and filling continuity give a
nonzero cycle $A$ maximizing
\[
 V(B)-c\M(B)^{(n+1)/n}\qquad(c>c_n).
\]
A constant rescaling makes its mass $m=\M(A)$ smaller than $s_n$.
This is the small-mass inequality that the rest of the proof contradicts.

The first step varies $A$ inward along geodesics from a fixed center
$y$. Write $r(x)=d(y,x)$, and let $Dr$ denote the differential of $r$
along the Euclidean current charts. Triangle comparison bounds the
deformation by a quadratic form in the time and chart variables.
Its determinant retains the factor $\sqrt{1-|Dr|^2}$ in the swept-mass
bound. In a smooth Euclidean model this is the transverse component of
the radial unit direction. The argument needs neither an ambient
normal bundle nor normality of a separate chart restriction.
Section~\ref{sec:radial} proves the resulting radial inequality.
For $n=2$, an inverse-square cutoff and Euclidean lower density
already force $m\ge4\pi=s_2$, establishing the base case.

For $n>2$, the induction hypothesis fills the boundaries of small
restrictions of $A$. Replacing each restriction by such a filling
gives an almost Euclidean small-set perimeter estimate. Intrinsic
coarea and rearrangement turn this into a critical Sobolev inequality
in Section~\ref{sec:sobolev}. The sharp coefficient is the classical
Aubin--Talenti coefficient~\cite{Aubin1976,Talenti1976}; its radial
form is proved here by the mass-transport method of
Cordero-Erausquin, Nazaret and Villani
\cite[Section~2]{CorderoNazaretVillani2004}.

To use this estimate variationally, we close the chart gradient in
$L^2(\mu)$, where $\mu=\|A\|$ is the mass measure. We prove compactness
by first projecting whole weighted currents and then isolating charts
in total variation. With $\beta=1/(n-2)$ and $p=2n/(n-2)$, the
almost sharp inequality and Br\'ezis--Lieb splitting
\cite{BrezisLieb1983} produce a nonnegative minimizer of
\[
 \frac{4\beta\int |Du|^2\,d\mu+n\int u^2\,d\mu}
 {\bigl(\int |u|^p\,d\mu\bigr)^{2/p}}
\]
on the completion for the norm
$(\|u\|_2^2+\int|Du|^2\,d\mu)^{1/2}$. Section~\ref{sec:minimizer} normalizes this
minimizer $v$ so that $0<W:=\int v^p\,d\mu<s_n$, and establishes
the moment estimates needed for its nonlinear tests.

The final comparison concerns the weighted chord
\[
 v(y)^\beta d(y,x)v(x)^\beta
\]
and the probability measure $d\nu=v^p\,d\mu/W$.
Radial variation and the minimizer equation control a scalar transform
of the chord distribution. Euclidean tangent density then bounds its
mean square by $2(W/s_n)^{2/n}<2$ for almost every center $y$.
The actual composite tests remain differentiable at $v=0$, even when
$\beta<1$; no Sobolev regularity of $v^\beta$ is presumed.
Section~\ref{sec:chords} supplies this upper bound.
Section~\ref{sec:conclusion} uses CAT(0) variance to show that the
average of the same squared chords is at least $2$. The contradiction
completes the induction and yields attained, compactly supported
integral fillings.
